% Final
\subsection{Scrimba}

% Final
Poslední analyzovanou aplikací je webová aplikace Scrimba\footnote{Dostupné na: \href{https://scrimba.com/}{https://scrimba.com/}}. Tato aplikace se od ostatních odlišuje zejména svým inovativním způsobem výuky programování, který kombinuje prostředí připomínající skutečné IDE a interaktivní výuková videa. Lekce se tak odehrávají přímo v editoru kódu, ve kterém je uživateli vysvětlována látka a zadávány úkoly. Tento přístup výuky může přinést řadu výhod oproti tradičním metodám výuky a proto jsem se rozhodl tuto aplikaci do analýzy zahrnout.

% Final
Platforma Scrimba obsahuje přes 80 kurzů na různá témata napříč softwarovým inženýrstvím \cite{ScrimbaCourses}. V době psaní tohoto textu jsou kurzy zaměřeny zejména na webové technologie a umělou inteligenci.

% Final
Na rozdíl od ostatních analyzovaných aplikací Scrimba nemá v současné době mobilní verzi aplikace. Proto pro analýzu využiji dostupnou webovou aplikaci Scrimba. 

% Final
\subsubsection{Vzhled a uživatelská zkušenost}

% Odstavec: Obecný dojem (komplexita/přehlednost UI, paleta barev, typografie)

% Final
Na první pohled působí Scrimba jako přehledná a minimalistická aplikace. Ve srovnání s konkurenční aplikací Mimo používá aplikace menší velikost písma a zobrazuje na stránkách více textu. Aplikace tak může působit komplexnějším dojmem. Co se palety barev týče, používá Scrimba monochromatickou paletu modré barvy. Odstíny modré jsou použity pro pozadí, primární akce, ikony i textové prvky. Pro upozornění a některé štítky aplikace výjimečně využívá žlutou nebo zelenou barvu. Napříč aplikací je použito systémové písmo, tedy písmo, které je nastaveno v závislosti na operačním systému uživatele. Na systému Windows je tak použito například písmo Segoe, na zařízeních macOS pak písmo San Francisco a na zařízeních Android písmo Roboto \cite{SystemFontWebSafeFont}.

% Odstavec: Rozložení prvků na obrazovce (lekce v kurzu, učení)
% Final
Prvky jsou na obrazovce rozloženy velice přehledně a logicky. Ve srovnání s ostatními aplikacemi Scrimba používá méně prostoru mezi jednotlivými prvky uživatelského rozhraní. V důsledku toho může aplikace působit komplexnějším dojmem.

% Final
Na hlavní stránce kurzu jsou jednotlivé sekce, moduly a lekce vizuálně uspořádány do stromové struktury, kterou lze postupně rozbalovat. Díky tomu může uživatel jednoduše získat přehled o celém kurzu a rozbalit části, které ho zajímají. Obsah kurzu je tak uspořádán velice přehledně.

% Odstavec: Navigace v aplikaci

% Final
Na rozdíl od ostatních analyzovaných aplikací Scrimba nevyužívá pro navigaci horní navigační lištu, ani žádnou patičku stránky. Pro navigaci je využívána výhradně boční navigační lišta, která je umístěna na levé straně obrazovky.

% Odstavec: Animace a efekty a mikro interakce

% Final
Ve srovnání s ostatními aplikacemi Scrimba výrazně více využívá animace jednotlivých prvků uživatelského rozhraní a animované přechody mezi stránkami. Na většině stránek aplikace jsou při načtení zobrazeny animace některých prvků rozhraní. Další výraznou animaci lze zpozorovat při přejetí kurzorem nad tlačítko. Při animaci je tlačítko zvýrazněno čarou, která ho postupně obklopí po jeho obvodu.

% Final
Animace jsou tak napříč aplikací používány často, nicméně minimalisticky a v přiměřeném množství. Díky tomu Scrimba může působit výrazně interaktivnějším a vizuálně poutavějším dojmem, než její konkurenti.

% Final
\subsubsection{Způsoby učení}

% Struktura kurzu (lineární, tematické moduly, career paths, kapitoly atd.)

% Final
Scrimba nabízí, podobně jako ostatní analyzované aplikace, kurzy a kariérní cesty. Kariérní cesty jsou takřka stejné jako kurzy a liší se pouze v tom, že obsahují více témat. Lekce v kurzech jsou uspořádány chronologicky a seskupeny do kapitol a podkapitol.

% Jak se spustí učení a jak vypadá study session

% Final
Unikátní funkcí této aplikace je samotné učení. Při spuštění učení je uživateli v prohlížeči zobrazen editor připomínající skutečné IDE, ve kterém je možné prohlížet soubory a editovat kód. V editoru zároveň uživatel vidí okno s výstupem daného kódu. Dále je uživateli zobrazena prezentace doplňující výklad dané látky. Všechna tato okna může uživatel v průběhu výkladu schovávat, zobrazovat a procházet dle libosti.

% Final
Při učení uživatel poslouchá audio s vysvětlením dané látky. V průběhu toho je mu buď zobrazena prezentace nebo editor, ve kterém může sledovat, jak lektor přepisuje kód a jak pohybuje svým kurzorem myši. Při výkladu může uživatel procházet libovolné soubory a přepisovat jejich obsah. Dále si může kdykoliv zobrazit okno s prezentací nebo výstupem kódu. Audio s výkladem látky může uživatel kdykoliv přetáčet zpět a nebo dopředu. V průběhu lekce jsou uživateli zadávány úkoly, které musí v editoru splnit. Zadání je provedeno lektorem v audio nahrávce. Po zadání úkolu je audio pozastaveno, dokud uživatel nesplní daný úkol. 

% Druhy cvičení
% Final
Na rozdíl od ostatních analyzovaných aplikací Scrimba nepoužívá předdefinované typy cvičení. Veškerý výklad látky i cvičení se odehrávají přímo v editoru. Cvičení typicky spočívají v úpravě existujícího kódu nebo napsání nového kódu. Celkově je učení velice interaktivní a připomíná kombinaci videí s výkladem a skutečného IDE. Tato forma učení může být pro uživatele velice přínosná, zejména kvůli názornějšímu vysvětlení látky a interaktivnímu způsobu učení.

% Final
\subsubsection{Způsoby motivace uživatele}

% Jaké formy gamifikace aplikace používá
% Final
Na rozdíl od ostatních aplikací Scrimba neobsahuje téměř žádné formy gamifikace. Uživatelé nesbírají žádné body nebo ocenění, a nemohou spolu v aplikaci interagovat. Motivující pro uživatele ovšem může být, že aplikace jasně ukazuje, kolik procent kurzu uživatel dokončil a kolik cvičení uživatel splnil. Všechny splněné lekce jsou také označeny výraznou zelenou barvou a ikonou zaškrtnutí. Tato vizuální indikace pokroku může pro některé uživatele sloužit jako zdroj motivace.

% Final
\subsubsection{Placená verze}

% Final
Scrimba na rozdíl od ostatních aplikací neobsahuje žádné reklamy a nabízí pouze jeden model předplatného.

% Free tier
% Final
Uživatelé bez předplatného mají k dispozici většinu kurzů zdarma. Některé kurzy a kariérní cesty jsou ovšem k dispozici pouze uživatelům s aktivním předplatným. V nich mají uživatelé bez předplatného možnost si vyzkoušet několik prvních lekcí.

% Premium tier
% Final
Uživatelé s aktivním předplatným mají kromě základních kurzů přístup k pokročilejším kurzům a kariérním cestám. Na konci každého kurzu si navíc mohou vygenerovat a stáhnout certifikát o dokončení daného kurzu.

% Final
\subsubsection{Silné a slabé stránky}

% Silné stránky
% Final
Silnou stránkou této aplikace je její unikátní přístup k učení. Uživatelé se učí přímo v editoru připomínající IDE a veškerá látka je vysvětlena lektorem a doplněna o prezentaci s obrázky, GIFy a diagramy. Učení je tak zábavné, interaktivní a může být pro uživatele více srozumitelné, zejména u složitější látky.

% Final
Další silnou stránkou aplikace je minimalistický design a rozložení obsahu. Aplikace se velice snadno používá a je velice jednoduché najít, co uživatel potřebuje. Důvodem může být také to, že aplikace obsahuje pouze funkce, které jsou klíčové pro výuku programování a nezahlcuje uživatele dodatečnými materiály a nástroji.

% Slabé stránky
% Final
Mírnou nevýhodou inovativního způsobu učení je volnost uživatele při učení. Uživatel může v průběhu výkladu svým počínáním například omylem vypnout okno s prezentací nebo výstupem kódu. Proto se může velice jednoduše stát, že lektor v audio nahrávce mluví například o diagramu v prezentaci, který se uživateli na obrazovce nezobrazuje, protože uživatel dříve kliknul na editor, což zamezilo zobrazení okna s prezentací. Pro uživatele tak může být místy velice matoucí, kde se zrovna lektor nachází v průběhu výkladu a může být těžké výklad sledovat.

% Final
Další potenciální nevýhodou může být nedostatek motivace uživatele. Aplikace neobsahuje téměř žádné funkce, které by se zaměřovaly na dlouhodobou motivaci uživatelů. Proto může být pro uživatele obtížné se dlouhodobě přimět k používání této aplikace.

% Final
Poslední slabou stránkou aplikace Scrimba je místy nefunkční editor kódu. Protože do kódu může zároveň zapisovat virtuální lektor a uživatel, několikrát se při mým průchodem aplikací stalo, že byl například kurzor posunutý o několik znaků vedle, nebo že psaní do kódu nefungovalo a přepisovaly se jiné znaky, než by uživatel očekával. Tento problém může být pro uživatele velice odrazující při používání aplikace, protože se pak aplikace jeví jako nefunkční.

% Final
\subsubsection{Další poznatky}

% Zajímavé funkce
% Final
V rámci analýzy jsem došel k pozoruhodnému poznatku. Na rozdíl od ostatních aplikací Scrimba nemá dedikovanou domovskou stránku, která by například uživateli prezentovala výhody aplikace a jiné marketingové materiály. Místo toho je uživateli na hlavní stránce zobrazen seznam všech kurzů a krátké video popisující, jak se aplikace používá. Uživatelské rozhraní nepřihlášeného a přihlášeného uživatele se tak téměř neliší. Uživatelé mohou prohlížet všechny kurzy bez přihlášení a mohou spustit až dvě libovolné lekce, než po nich bude aplikace vyžadovat přihlášení. Díky tomu mohou uživatelé zjistit, zda jim vyhovuje formát výuky, než se do aplikace zaregistrují.

% Final
\subsubsection{Celkové hodnocení}

% Final
Souhrnné hodnocení aplikace Scrimba je uvedeno v tabulce \ref{table:competitor-rating-scrimba}.

% Final
\begin{table}[ht]
    \centering
    \begin{tabular}{@{} l c @{}}
        \toprule
        Oblast hodnocení & Hodnocení \\
        \midrule
        Uživatelská zkušenost & 8/10 \\
        Jednoduchost učení & 5/10 \\
        Podpora opakování & 3/10 \\
        Motivace uživatelů & 3/10 \\
        Interaktivita cvičení & 9/10 \\
        \bottomrule
    \end{tabular}
    \captionsetup{justification=centering}
    \caption{Hodnocení aplikace Scrimba}
    \label{table:competitor-rating-scrimba}
\end{table}
